\chapter{Introduction}\label{chap:intro}
Le groupe CHU souhaite un entrepôt de données décisionnel (EDD) unifiant quatre sources hétérogènes : une base PostgreSQL opérationnelle (S1), des fichiers CSV de l'annuaire santé (S2), des fichiers de satisfaction HAS accessibles par FTP (S3) et le fichier des décès de l'INSEE (S4).

Ce livrable 1 présente les quatre composantes demandées : le modèle conceptuel et logique du schéma décisionnel (chapitre~\ref{chap:modele}), l'architecture cible Bronze/Silver/Gold avec les technologies retenues (chapitre~\ref{chap:archi}), les spécifications fonctionnelles des jobs ETL (chapitre~\ref{chap:etl}) et le plan de gouvernance des données (chapitre~\ref{chap:gouv}). Le dictionnaire de données complet est en annexe~\ref{sec:dictionnaire}.
